\section{The equation and its polynomial cylinder}\label{sec:construction}

We first prove that the displayed equation defines an integral fourfold,
then construct the five polynomial coordinates on its cylinder. The
localization at $p$ explains the equation; the stabilization argument
will produce coordinates in the original ring, including the fiber
$p=0$.

The construction realizes the residual-coordinate method of
El Kahoui--Ouali~\cite[Lemma~3.1 and Theorem~1.1]{ElKahouiOuali},
generalized by Dutta--Lahiri~\cite[Lemma~3.3 and Theorem~3.4]{DuttaLahiri}.
The coordinate criterion of Maubach~\cite[Theorem~4.5]{Maubach}
applies to $H+p^3w$: the only root of $p^3$ is $p=0$, where
the coordinate change below makes $H$ a coordinate. We retain the
explicit formulas so that the resulting isomorphism is available
without an existence argument.

Recall the polynomial ring and its elements
\begin{equation}\label{con:setup}
\begin{gathered}
 P=\kk[p,s,u,F,J],\qquad x=s^2+u^3+p^2F,\\
 y=s+x(x-u^3),\qquad z=sx+p^2J,\\
 H=x^2F-(1+2sx)J-p^2J^2-pu,\qquad A=P/(H).
\end{gathered}
\end{equation}
Thus \(H\) is a specified polynomial in the five independent variables
\(p,s,u,F,J\). We will use \(w\) for the polynomial indeterminate adjoined
to \(A\). All maps below are \(\kk\)-algebra maps.

\subsection{Localization and integrality}

The definitions give the identity
\begin{equation}\label{con:identity}
 xy-z(z+1)=p^2(H+pu).
\end{equation}
Indeed, substituting \(y\) and \(z\) into the left-hand side and cancelling
the terms \(sx\) gives
\[
 x^2(x-u^3-s^2)-p^2(1+2sx)J-p^4J^2,
\]
which is the right-hand side because \(x-u^3-s^2=p^2F\).

\begin{lemma}\label{con:domain}
The algebra \(A\) is an integral domain of dimension four, \(p\ne0\)
in \(A\), and
\[
 A[p^{-1}]=\kk[p,p^{-1},x,y,z]
\]
is a Laurent polynomial ring in the displayed generators.
\end{lemma}

\begin{proof}
Over \(\kk[p,p^{-1}]\), the elements \(x,y,z,u\) are polynomial
coordinates on \(P[p^{-1}]\). The inverse substitutions are
\begin{equation}\label{con:localized-coordinates}
 s=y-x(x-u^3),\qquad
 F=p^{-2}(x-s^2-u^3),\qquad J=p^{-2}(z-sx).
\end{equation}
Substitution in either direction gives the identity, so in particular
the new coordinates are algebraically independent. In these coordinates,
\Cref{con:identity} reads
\begin{equation}\label{con:localized-h}
 H=p^{-2}\bigl(xy-z(z+1)\bigr)-pu.
\end{equation}
The coefficient of \(u\) is a unit. Replacing \(u\) by \(H\) is therefore
another polynomial coordinate change over \(\kk[p,p^{-1}]\). Consequently
\((H)\) is prime in \(P[p^{-1}]\), and its quotient is the Laurent
polynomial ring asserted in the lemma, with
\(u=p^{-3}(xy-z(z+1))\).

Set \(x_0=s^2+u^3\). Reduction modulo \(p\) gives
\[
 H\bmod p=x_0^2F-(1+2sx_0)J\ne0,
\]
so \(p\) does not divide \(H\) in \(P\). To descend irreducibility from
the localization, factor \(H\) in the unique factorization domain \(P\).
Its image in \(P[p^{-1}]\) is irreducible, so all but one of its
irreducible factors must become units there. An irreducible element of
\(P\) that becomes a unit upon inverting \(p\) divides a power of \(p\),
and hence is associated to \(p\). Nondivisibility by \(p\) excludes such
factors. Thus \(H\) is irreducible in \(P\), and is prime because \(P\)
is a unique factorization domain. This proves that \(A\) is a domain.
Its localization computed above is nonzero, which also proves that
\(p\ne0\) in \(A\). This localization has fraction field
\(\kk(p,x,y,z)\), of transcendence degree four over \(\kk\).
The same is true of \(A\), so the dimension formula for finitely
generated domains over a field~\cite[Tag~00P0]{Stacks} gives \(\dim A=4\).
\end{proof}

\subsection{An explicit stabilization}

The main calculation has two stages. An exponential automorphism
replaces $H$ by $H+p^3w$. A linear coordinate change in $F,J$ then
makes this new equation linear modulo $p^2$, which allows a polynomial
root modulo $p^3$ and a fifth coordinate to be written down explicitly.

\begin{proposition}\label{prop:stabilization}
For the algebra \(A\) in \Cref{con:setup}, there is an algebraic
\(\kk\)-isomorphism
\[
 A[w]\cong\kk\poly{5}.
\]
\end{proposition}

\begin{proof}
We first change the equation by a polynomial automorphism, and then
exhibit five polynomial coordinates on the resulting quotient.

\smallskip
\noindent\emph{Changing the equation.}
On \(P[p^{-1}]=\kk[p,p^{-1},x,y,z,u]\), define
\(\Delta=-p^2\partial/\partial u\), with \(p,x,y,z\) fixed.
This derivation is locally nilpotent: every polynomial is killed by a
sufficiently high iterate of \(\Delta\). Applying it to
\Cref{con:localized-coordinates} gives
\begin{equation}\label{con:derivation}
\begin{gathered}
 \Delta p=0,\qquad \Delta u=-p^2,\qquad
 \Delta s=-3p^2xu^2,\\
 \Delta F=(6sx+3)u^2,\qquad \Delta J=3x^2u^2.
\end{gathered}
\end{equation}
These values belong to \(P\), so \(\Delta\) restricts to a locally
nilpotent derivation of \(P\). Extend it by \(\Delta w=0\). The formula
\[
 \Phi(f)=\exp(w\Delta)(f)
       =\sum_{i\ge0}\frac{w^i}{i!}\Delta^i(f)
       \qquad(f\in P[w])
\]
is a finite sum for every \(f\). The iterated Leibniz rule makes \(\Phi\)
an algebra homomorphism, and the exponential identity gives its inverse
\(\exp(-w\Delta)\). By \Cref{con:localized-h},
\(\Delta H=p^3\) and \(\Delta^2H=0\). Hence
\begin{equation}\label{con:cylinder-automorphism}
 \Phi(H)=H+p^3w,\qquad
 A[w]\xrightarrow[\ \Phi\ ]{\ \sim\ }
 T:=P[w]/(H+p^3w).
\end{equation}
In particular, \(T\) is a domain by \Cref{con:domain}, and its element
\(p\) is nonzero since \(\Phi\) fixes \(p\).

\smallskip
\noindent\emph{Coordinates adapted to the fiber \(p=0\).}
With \(x_0=s^2+u^3\), put
\begin{equation}\label{con:linear-coordinates}
 L=x_0^2F-(1+2sx_0)J,\qquad
 M=(1-2sx_0)F+4s^2J.
\end{equation}
The determinant of this linear change in \((F,J)\), over
\(\kk[p,s,u]\), is
\[
 4s^2x_0^2+(1+2sx_0)(1-2sx_0)=1.
\]
Its inverse is
\begin{equation}\label{con:linear-inverse}
 F=4s^2L+(1+2sx_0)M,\qquad
 J=-(1-2sx_0)L+x_0^2M.
\end{equation}
It follows that \(p,s,u,L,M\) are algebraically independent and that
\[
 P=B[L],\qquad B:=\kk[p,s,u,M]\cong\kk\poly{4}.
\]
For the rest of the proof, regard \(F,J,H\) as polynomials in \(L\)
over \(B\), using \Cref{con:linear-inverse}. Expansion only in powers
of \(p\) gives
\begin{equation}\label{con:h-expansion}
 H(L)=L-pu+p^2Q(L)+p^4F(L)^3,
 \qquad Q(L)=2x_0F(L)^2-2sF(L)J(L)-J(L)^2.
\end{equation}

At this point the transfer theorem of
Edo--V\'en\'ereau~\cite[Theorem~7]{EdoVenereau} also applies, with
coefficient ring $B$ and parameter $p^3$: every coefficient of
$H(L)-L$ is divisible by $p$, hence nilpotent modulo $p^3B$.
The following calculation gives a polynomial inverse in this instance.

\smallskip
\noindent\emph{A polynomial root modulo \(p^3\).}
Define elements and a polynomial over \(B\) by
\begin{equation}\label{con:root}
 \begin{gathered}
 C=Q(0)=M^2(2x_0+6sx_0^2+4s^2x_0^3-x_0^4),\\
 Q_1(L)=\frac{Q(L)-C}{L}\in B[L],\qquad
 L_*=pu-p^2C\in B.
 \end{gathered}
\end{equation}
Since \(Q(L_*)-C=L_*Q_1(L_*)\), \Cref{con:h-expansion} gives the exact
identity
\begin{equation}\label{con:root-divisibility}
 H(L_*)=p^3\bigl((u-pC)Q_1(L_*)+pF(L_*)^3\bigr)\in p^3B.
\end{equation}

\smallskip
\noindent\emph{A polynomial representative for the fifth coordinate.}
The quotient $(L-L_*)/p^3$ suggests the remaining coordinate after
localization. The next identity provides a polynomial representative
for it before $p$ is inverted.
In \(B[L,w]\), put
\begin{equation}\label{con:e-polynomial}
 h=u-pQ(L)-p^3F(L)^3-p^2w,
 \qquad e_0=-w-pF(L)^3-Q_1(L)h.
\end{equation}
Then \(H(L)+p^3w=L-ph\). Using \(Q(L)-C=LQ_1(L)\), we obtain the
polynomial identity
\begin{equation}\label{con:e-certificate}
 p^3e_0-(L-L_*)
   =(p^2Q_1(L)-1)\bigl(H(L)+p^3w\bigr).
\end{equation}
Thus the image \(e\) of \(e_0\) in \(T\) satisfies
\begin{equation}\label{con:e-relation}
 p^3e=L-L_*.
\end{equation}
In particular, \(e\) belongs to the unlocalized ring \(T\).

\smallskip
\noindent\emph{Generation and algebraic independence.}
Inverting \(p\) in the presentation
\(T=B[L,w]/(H(L)+p^3w)\) allows us to eliminate \(w\), giving the
isomorphism
\begin{equation}\label{con:t-localized}
 T[p^{-1}]\xrightarrow{\ \sim\ }B[p^{-1}][L],
 \qquad w\longmapsto-p^{-3}H(L),
\end{equation}
which is the identity on \(B\) and \(L\). The composite
\(B\to T\to T[p^{-1}]\) is therefore the natural injection
\(B\hookrightarrow B[p^{-1}][L]\). This proves that the displayed copy
of \(B\) embeds in \(T\). Moreover, \Cref{con:e-relation} identifies
\(e\) in this localization with \(p^{-3}(L-L_*)\), an affine linear
coordinate in the independent variable \(L\). Consequently \(e\) is
transcendental over \(B[p^{-1}]\), and hence over \(B\).

It remains to prove that \(B\) and \(e\) generate \(T\).
By \Cref{con:e-relation}, \(L=L_*+p^3e\). The polynomial difference
\(H(L_*+p^3Z)-H(L_*)\) is divisible by \(p^3\) in \(B[Z]\).
Together with \Cref{con:root-divisibility}, this shows that
\[
 W(Z):=-\frac{H(L_*+p^3Z)}{p^3}\in B[Z].
\]
The defining relation of \(T\) now gives
\(p^3(w-W(e))=0\). Since \(T\) is a domain and \(p\ne0\), it follows
that \(w=W(e)\). Thus \(L,w\in B[e]\), which proves
\begin{equation}\label{con:polynomial-cylinder}
 T=B[e]=\kk[p,s,u,M,e]\cong\kk\poly{5}.
\end{equation}

For completeness, the five independent coordinates in \(A[w]\) itself
are the classes of
\[
 \Phi^{-1}(p),\quad\Phi^{-1}(s),\quad\Phi^{-1}(u),\quad
 \Phi^{-1}(M),\quad\Phi^{-1}(e_0).
\]
These are polynomials by the finite exponential formula. They are
obtained by transporting the five generators in
\Cref{con:polynomial-cylinder} through the inverse of
\Cref{con:cylinder-automorphism}, which proves the proposition.
\end{proof}
